Ce projet étudie le pilotage conjoint de la maintenance et du stock de pièces de rechange d'une flotte de 100 machines réparties en deux familles, l'une vieillissant à l'usage, l'autre au calendrier. Un simulateur à événements discrets reproduit pannes, maintenances, approvisionnements, défauts de fabrication et télémétrie ; calibré sur le jeu Microsoft Azure Predictive Maintenance, il retrouve le nombre de pannes de chaque composant sur l'année 2015 à moins de 10~\% près. Des réseaux de survie récurrents et convolutifs estiment le risque de panne de chaque composant à partir de 48 heures de signaux ; ils dépassent LightGBM sur données simulées et, à 7 jours, sur données réelles. Une politique prédictive fondée sur ces réseaux réduit le coût annuel de 2,1~\% par rapport à une politique à seuils fixes optimisés, écart significatif sur 30 années simulées, mais une maintenance conditionnelle simple fait mieux encore, avec 9,6~\% d'économie. Les agents d'apprentissage par renforcement profond (PPO avec masquage d'actions, DQN, Branching Dueling Q-Network) n'atteignent pas la politique à seuils optimisés : le meilleur, qui règle chaque jour les seuils de maintenance et de stock, s'en approche à 1,1~\% près, tandis que les agents qui décident machine par machine restent loin derrière. Un modèle de substitution ($R^2$ de 0,98 sur le coût) permet l'analyse de sensibilité de Sobol, et une application web restitue l'ensemble.
